\section{When indexes are built and when queries use them}
\label{sec:index-runtime-selection}
\label{sec:index-registry}

Index maintenance and index reading are separate decisions. Saving an index
does not force every query to use it, and selecting an index for reading does
not build it. The policies below concern BioseqDB's sequence indexes, not
PostgreSQL's native indexes on table columns.

\subsection{Maintenance policies}
\label{sec:index-update-policy}

Each registered index has a maintenance policy for each sequence column.
\path{bioseq.set_index_policy} changes that policy for all current sequence
columns in the named table. Table~\ref{tab:index-update-policy} describes the
two choices, and Table~\ref{tab:index-policy-defaults} gives their defaults.

\begin{table}[htbp]
\centering
\small
\caption{Maintenance policies selected with \texttt{bioseq.set\_index\_policy}.}
\label{tab:index-update-policy}
\begin{tabularx}{\linewidth}{@{}p{.31\linewidth}X@{}}
\toprule
\textbf{Policy} & \textbf{Meaning}\\
\midrule
\texttt{UPDATE\_ON\_INSERT} & Request automatic maintenance for newly stored
sequences when the global insertion-maintenance setting permits it. Ordinary
indexes are built for the new sequence, without rebuilding older rows.\\
\texttt{UPDATE\_MANUALLY} & Skip automatic insertion maintenance for this
index. Build or rebuild it with \texttt{bioseq.update\_index}. Existing index
contents remain available to eligible readers.\\
\bottomrule
\end{tabularx}
\end{table}

\begin{table}[htbp]
\centering
\small
\caption{Default policies before a table has explicit overrides.}
\label{tab:index-policy-defaults}
\begin{tabularx}{\linewidth}{@{}p{.34\linewidth}p{.30\linewidth}X@{}}
\toprule
\textbf{Index} & \textbf{Maintenance policy} & \textbf{Default read role}\\
\midrule
\path{minimizers_per_fragment} & \texttt{UPDATE\_ON\_INSERT} & Active minimizer index\\
\path{minimizers_per_sequence} & \texttt{UPDATE\_MANUALLY} & Not selected\\
\path{symbol_counts} & \texttt{UPDATE\_ON\_INSERT} & Used by supported symbol-count operations\\
\bottomrule
\end{tabularx}
\end{table}

There is also a global setting,
\path{bioseq.index.insert_maintenance.enabled}, whose current default is
\texttt{off}. Thus a default \texttt{INSERT} saves the sequence but builds
none of these indexes. In particular, the MPF policy's name does not imply
that it is being maintained while this setting is off.

The following example deliberately selects ordinary MPF maintenance.
Disabling grouped construction is necessary because the ordinary insertion
hook cannot publish a grouped minimizer generation. I also select the scalar
algorithm because the default bounded ntHash requires the grouped format to
record its hash identity. The table uses
\texttt{bioseq.seq('dna4',32)} and starts without a grouped MPF index.

\noindent\begin{minipage}{\linewidth}
\begin{lstlisting}[style=bioseq-index-sql,
  caption={Enable ordinary MPF maintenance for an inserted sequence.},
  label={lst:index-policy}]
SET bioseq.aligner.grouped_soa_generation.enabled = off;
SET bioseq.aligner.minimizer.algorithm = 'scalar';
SET bioseq.index.insert_maintenance.enabled = on;
SELECT bioseq.set_index_policy(
  'reads'::regclass, 'minimizers_per_fragment', 'UPDATE_ON_INSERT');

INSERT INTO reads (id, sample_id, seq)
VALUES (101, 7, repeat('ACGT', 2000)::bioseq.seq('dna4', 32));
\end{lstlisting}
\end{minipage}

The switches in Listing~\ref{lst:index-policy} are alternatives to the
grouped workflow in Section~\ref{sec:index-mpf5-sidecar}. They are not
additional steps to run after that example. They also permit maintenance of
other indexes with an on-insert policy, including symbol counts.
An \texttt{UPDATE} writes a replacement sequence row and applies insertion
maintenance to that replacement. Changing the policy does not backfill
older rows, erase existing indexes, or choose the active read index.

When grouped construction and insertion maintenance are both enabled,
ordinary \texttt{INSERT} rejects an enabled minimizer index with on-insert
policy. For this workflow, the user must leave
insertion maintenance off, commit the source rows and explicitly build their
index. This prevents an ordinary relation-only write from leaving a grouped
reader file inconsistent, as explained in
Section~\ref{sec:index-file-consistency}.

\texttt{bioseq.load} has a separate batched path. When insertion maintenance
is enabled, it first finishes the streamed source writes, then calls the
explicit builder for indexes with \texttt{UPDATE\_ON\_INSERT}. This permits
grouped construction after loading within that command. It is not the
per-row \texttt{INSERT} hook, and the grouped builder's usual restrictions
still apply.

For example, this policy keeps MPS outside automatic load maintenance:

\noindent\begin{minipage}{\linewidth}
\begin{lstlisting}[style=bioseq-index-sql,
  caption={Keep MPS construction separate from loading.}]
SELECT bioseq.set_index_policy('reference_sequences'::regclass,
  'minimizers_per_sequence', 'UPDATE_MANUALLY');
\end{lstlisting}
\end{minipage}

MPS already has this policy by default. The explicit call is useful when a
table previously selected automatic MPS maintenance. It changes neither MPF
nor symbol-count policy.

After loading and committing the source rows, the next listing builds MPS.
It assumes the current defaults, with generation construction enabled and both
asynchronous-persistence switches off. The target starts without an older
ordinary MPS relation. Changing a construction setting does not convert an
existing ordinary relation into a grouped one.

\noindent\begin{minipage}{\linewidth}
\begin{lstlisting}[style=bioseq-index-sql,
  caption={Build MPS explicitly after loading and committing the source.},
  label={lst:index-build-mps}]
BEGIN;
SELECT bioseq.update_index('reference_sequences'::regclass,
                          'minimizers_per_sequence');
COMMIT;
\end{lstlisting}
\end{minipage}

The builder uses the configured algorithm, $k$ and $w$ for a supported
compiled profile. The explicit rebuild runs in the caller's transaction and
does not depend on the insertion-maintenance gate. Supported CPU MPS profiles
use generation storage grouped by hash across the sequence. Other supported
MPS routes use the grouped arrays described earlier. Neither creates the MPF
reader file.

\subsection{Reading policies}
\label{sec:index-read-policy}

A search has a per-call \texttt{index\_read\_requirement}, supplied through
\texttt{aligner\_options}. Its default is \texttt{allow}, encoded as 0.
It is not a GUC. A GUC is a PostgreSQL configuration setting changed with
\texttt{SET}, such as \path{bioseq.index.sequence_read.enabled}.

\begin{table}[htbp]
\centering
\small
\caption{The search's \texttt{index\_read\_requirement} option.}
\label{tab:index-read-requirement}
\begin{tabularx}{\linewidth}{@{}p{.31\linewidth}X@{}}
\toprule
\textbf{Value} & \textbf{What the search requires}\\
\midrule
0: \texttt{allow} (default) & Permit a supported compatible stored read when
\path{bioseq.index.sequence_read.enabled} is on. The executor can instead
compute a temporary lookup. This value does not select the grouped MPF CPU
reader described in Section~\ref{sec:index-mpf5-sidecar}.\\
1: \texttt{require\_absent} & Do not read a persisted index or reuse its
assembled lookup. Compute the lookup for this execution. The name refers
to execution without persisted reads, not to an assertion that no index
exists on disk.\\
2: \path{require_compatible_persisted} & Require a compatible saved
generation through a supported reader. A missing or incompatible generation
is an error. Recomputing it during the query cannot satisfy this request.
This requirement overrides the optional-read GUC.\\
\bottomrule
\end{tabularx}
\end{table}

The optional-read GUC defaults to \texttt{on}. Setting it to \texttt{off}
prevents optional stored reads, but does not prevent index construction.
Conversely, setting it to \texttt{on} does not create missing index data.
Requiring a temporary lookup does not prohibit separately permitted
background persistence after the query.

Compatibility means more than finding a file. The saved generation must
belong to the visible source sequence and use the minimizer profile expected
by the query. An index built with one hash algorithm cannot supply the
positions for another algorithm's hash values. Changing $k$ or $w$ can also
change which minimizers exist. The reader checks these identities before
using the saved data.

The executor must support the requested read as well. A complete stored
target in an eligible two-input \texttt{BioseqBatchJoin} can use the grouped
MPF file with requirement 2. An ordinary single-table \texttt{BioseqScan}
currently rejects that requirement. It does not become a batch join because
a saved index exists.

\texttt{bioseq.set\_active\_read\_index(table, index\_name)} chooses which
logical minimizer index the table's sequence columns should use. It does not
build that index. Existing maintenance policies are retained. If the call
creates a new policy row, that row starts with manual maintenance, so an
explicit selection can differ from the implicit defaults in
Table~\ref{tab:index-policy-defaults}.

\warn{Alpha software and restrictive memory checks}{
\label{warn:index-query-minimizers}
Please note that BioseqDB is still alpha software and has numerous issues.
One problem is that its memory checks are too restrictive for small queries.
The current policy keeps at least 2~GiB of system memory available, so even
an allocation of a few hundred bytes can be refused when other applications
use too much RAM. You should leave more than 2~GiB available, plus space for
the query itself. On macOS, the estimate counts free and inactive pages.
There is also a process-memory limit of 80\% of physical RAM.

This can produce the misleading error
\texttt{transient query execution stats disagree with its generated backend}.
The query may simply have failed to allocate its minimizer buffer. It does
not mean that the target index is missing or that query minimizers rejected
that index. The saved index describes the target sequence. BioseqDB must
still compute minimizers for the query sequence to find matching hashes.

If this happens under memory pressure, close other applications and retry.
Rebuilding the target index does not fix the memory shortage. The error
handling needs to report allocation failures separately from actual
minimizer-statistics errors. If the message remains when enough memory is
available, that requires further investigation.
}

\subsection{Construction, deferred work and other settings}
\label{sec:index-build-settings}

Table~\ref{tab:index-main-settings} collects the switches that decide whether
index work is performed. They supplement the per-index and per-query policies.
The defaults here are those of the current implementation.

\begin{table}[htbp]
\centering
\small
\caption{BioseqDB index construction and reading switches.}
\label{tab:index-main-settings}
\begin{tabularx}{\linewidth}{@{}p{.39\linewidth}p{.08\linewidth}X@{}}
\toprule
\textbf{Setting} & \textbf{Default} & \textbf{Effect}\\
\midrule
\path{bioseq.index.insert_maintenance.enabled} & off & Global gate for automatic
maintenance, subject to each index's policy and the grouped-insertion
restriction described above.\\
\path{bioseq.index.sequence_read.enabled} & on & Allow optional stored reads.
Per-call requirements 1 and 2 override this choice.\\
\path{bioseq.aligner.grouped_soa_generation.enabled} & on & Select generation-based
minimizer builds. MPF publishes grouped arrays and a reader file. Supported
CPU MPS profiles publish data grouped by hash across the sequence.
Changing this setting alone builds or converts nothing.\\
\path{bioseq.index.async_persistence.enabled} & off & Permit eligible deferred
persistence work. Explicit \texttt{update\_index} remains synchronous.\\
\path{bioseq.index.recompute_async_persistence.enabled} & off & Additionally
permit a recomputing query to request deferred persistence. It needs the
preceding switch and does not guarantee a saved result.\\
\bottomrule
\end{tabularx}
\end{table}

The deferred implementation is limited. Its surviving query-side submission
path handles ordinary MPS, with $k=15,w=10$ and an accepted backend.
It excludes the default bounded ntHash algorithm. It also requires an
eligible enabled index with on-insert policy and an available worker queue.
Its policy selection prefers MPF before MPS, independently of the active
read index, but the submission path rejects MPF. Thus enabling both switches
alone is not enough to arrange an MPS rebuild.

Accepted requests retain source identifiers and parameters. They are
submitted after transaction commit and discarded on abort. A worker opens
its own transaction and checks the source and maintenance policy again.
It does not inherit the query's snapshot or relation pointers. This machinery
does not maintain grouped MPF reader files. The explicit grouped builder
rejects either asynchronous-persistence switch being on.

There are additional controls for representation and resource use, listed
in Table~\ref{tab:index-other-settings}. None introduces another read-policy
value or a separate switch for creating the MPF reader file.

\begin{table}[htbp]
\centering
\small
\caption{Other controls that affect BioseqDB sequence indexes.}
\label{tab:index-other-settings}
\begin{tabularx}{\linewidth}{@{}p{.40\linewidth}X@{}}
\toprule
\textbf{Setting and default} & \textbf{Role}\\
\midrule
\path{bioseq.aligner.minimizer.algorithm}\newline
\texttt{bounded\_nthash\_simd} & Selects minimizer generation and its hash
identity. Construction and reading must agree on a compatible profile.\\
\path{bioseq.aligner.minimizer.kmer_size}: 15\newline
\path{bioseq.aligner.minimizer.window_size}: 10 & Select $k$ and $w$.
The grouped builder validates them against compiled profiles.\\
\path{bioseq.aligner.fragment_minimizer.lookup_route}\newline
\path{direct_compact_hash_lookup} & Use the direct section for supported
32-bit MPF hashes. The alternative \texttt{control\_block\_scan} searches
fragment controls. Wider hashes use controls even under the direct default.\\
\path{bioseq.index.memory_cap_mb}: 0 & Additional resident-index memory cap.
Zero uses the common runtime budget. It does not mean unlimited memory.\\
\path{bioseq.index_build.run_buffer_bytes}: 0 & Internal superuser setting
for external-sort run size. Zero chooses a size from the memory budget.\\
\path{bioseq.index.sequence_minimizer_sort.algorithm}\newline
\texttt{std\_sort} & Select the ordinary MPS sorting implementation.
The other choices are \texttt{simd\_u128\_vqsort} and
\texttt{simd\_u128\_vec\_stream\_1}, \texttt{\_2}, or \texttt{\_4}.\\
\path{bioseq.aligner.vulkan.index_build.enabled}: off & Permit the GPU
index-build route when Vulkan execution and the selected algorithm support it.\\
\path{bioseq.aligner.vulkan.index_cache.enabled}: on & Retain compatible
GPU index data for reuse. A search may override it through
\texttt{aligner\_options(vulkan\_index\_cache\_enabled => ...)}.\\
\bottomrule
\end{tabularx}
\end{table}

Algorithm and occurrence-filter settings still affect the search itself.
For example, the query's occurrence limit must fit the direct file's saved
limit from Section~\ref{sec:index-direct-lookup}. CPU or GPU selection and
scheduler resources also determine which implementation can run. Those are
execution choices, not additional policies for maintaining stored indexes.

\subsection{Reusing bytes already in memory}
\label{sec:index-residency}

Persistent storage avoids rebuilding an index after every query. Memory
caches avoid some of the reads and preparation within a backend's lifetime.
They hold different objects:
\begin{itemize}
  \item \textbf{Serialized chunks.} Ordinary relation readers have a bounded
    cache of index bytes. A reader may still copy a cached chunk into its
    own buffer.
  \item \textbf{Prepared CPU lookups.} Ordinary MPF and MPS readers can retain
    assembled objects identified by source, index, algorithm and interval.
  \item \textbf{Mapped file pages.} The operating system caches pages touched
    through the reader-file mapping. The grouped route does not fill the
    ordinary serialized-chunk cache.
  \item \textbf{GPU arrays.} Device data can remain available while a batch
    or a compatible cache entry owns it.
\end{itemize}

Evicting a cache entry removes the cache's reference. An active reader can
still own the object, so a cache limit is not a bound on the backend's total
memory. Query buffers and other readers can overlap with it. Similarly,
a file's mapped address range is not the amount of RAM currently occupied
by its pages.

Abort cleanup discards ordinary chunk-cache entries written by the failed
transaction. Source and index invalidation discard affected cached objects.
Reusing a prepared lookup still requires matching identities. Releasing
these memory resources does not remove the stored index or reader file.

\section{What the database integration adds}
\label{sec:index-storage-conclusions}
\label{sec:index-operational}

BioseqDB keeps minimizer indexes alongside versioned sequence storage.
Ordinary maintenance can index newly inserted sequences, while grouped
construction publishes complete generations for selective readers. The MPF
reader file gives workers direct access to those saved arrays. These choices
avoid repeated work, but require explicit rules for construction, compatible
reads and source changes. Recovery of unlogged data and removal of retired
files remain limitations of the current implementation.
